\section{Clause omission and root tests}
\label{sec:deletion}

We now turn the lifetime estimate into a bound for the effect of
deleting one clause.  Its variables provide tests of the assignments
that remain after deletion.  We expose the incidence pattern through
the whole cap, so that the number of tests is known, while keeping
both their contents and the future ordinary clauses independent.
The tests bound the probability of a deletion event; the ordinary
clauses bound its remaining duration.

Fix an integer \(k\ge3\), a real \(B>0\), and \(n\ge2k\).  Put
\(L=\lfloor Bn\rfloor\), and for now suppose \(L\ge1\).
Let \(C_1,C_2,\ldots\) be independent uniform proper \(k\)-clauses
on \(n\) variables, let \(H\) be the first unsatisfiable prefix index,
and put \(T=\min\{H,L\}\).  For \(1\le i\le L\), define
\[
 B_m^{(i)}=\bigwedge_{\substack{1\le j\le m\\j\ne i}}C_j,\qquad
 T^{(i)}=\min\bigl\{\inf\{m\ge0:B_m^{(i)}\notin\sat\},L\bigr\},
 \qquad D_i=T^{(i)}-T.
\]
The infimum is over integers and \(\inf\varnothing=\infty\).
Deleting \(C_i\) retains all other clause indices.  It can only delay
unsatisfiability, so \(D_i\ge0\); moreover \(T^{(i)}\) depends only
on the first \(L\) clauses other than \(C_i\).

\begin{lemma}[Coordinate omission]\label{lem:coordinate-omission}
For the independent clause process above,
\begin{equation}\label{eq:coordinate-omission}
 \Var(T)\le\sum_{i=1}^{L}\EE D_i^2.
\end{equation}
\end{lemma}

\begin{proof}
This is the coordinate-omission form of the Efron--Stein
method~\cite{EfronStein}, recorded in
\cite[equation~(3.6)]{BBLM}.  Here is a direct martingale proof.
Let \(\mathcal H_j=\sigma(C_1,\ldots,C_j)\), with \(\mathcal H_0\)
trivial, and write
\[
 M_i=\EE[T\mid\mathcal H_i]-\EE[T\mid\mathcal H_{i-1}].
\]
Since \(T^{(i)}\) omits the independent coordinate \(C_i\),
\(\EE[T^{(i)}\mid\mathcal H_i]
=\EE[T^{(i)}\mid\mathcal H_{i-1}]\).  Thus, with
\(X_i=T-T^{(i)}\),
\[
 M_i=\EE[X_i\mid\mathcal H_i]-\EE[X_i\mid\mathcal H_{i-1}].
\]
The tower property gives
\[
 \EE M_i^2
 =\EE\bigl(\EE[X_i\mid\mathcal H_i]\bigr)^2
  -\EE\bigl(\EE[X_i\mid\mathcal H_{i-1}]\bigr)^2
 \le\EE X_i^2=\EE D_i^2,
\]
where the inequality uses conditional Jensen.  The martingale
differences are orthogonal and \(T\) is \(\mathcal H_L\)-measurable,
so summing proves the claim.
\end{proof}

\subsection{What the two filtrations reveal}

Fix \(i\) and condition on a value of \(C_i\).  Until the end of
Section~\ref{sec:occupation}, probabilities and expectations refer to
this conditional law.  The bounds will be uniform in the fixed clause.
Its set \(R\) of \(k\) variables will be called the roots, and
\(t\in\{0,1\}^{R}\) is their unique assignment falsifying \(C_i\).
There are \(u=n-k\ge k\) remaining variables.  Write \(B_m=B_m^{(i)}\).

For \(i\le m<L\), define the deletion event
\begin{equation}\label{eq:deletion-events}
 A_m=\{B_m\in\sat,\ B_m\wedge C_i\notin\sat\}
     =\{T\le m<T^{(i)}\}.
\end{equation}
On this event every solution of \(B_m\) has root values \(t\).
The two prefix processes agree before \(i\), and hence
\begin{equation}\label{eq:delay-identity}
 D_i=\sum_{m=i}^{L-1}\mathbf1_{A_m}.
\end{equation}
In particular \(D_L=0\): the displayed sum is then empty.  The upper
endpoint is \(L-1\) because both failure times have already been
stopped at \(L\).

For \(i\le m\le L\), define \(S_m\subseteq\{0,1\}^{u}\) to be the
nonroot restrictions of solutions of \(B_m\) with root values \(t\).
The sets are defined on every outcome and decrease with \(m\); for
\(m<L\), \(A_m\) implies that \(S_m\) is nonempty.  We need to reveal
this set while retaining independent randomness to test whether one
of its assignments allows a root to be flipped.

The \emph{root type} of a clause records the subset of its variables
in \(R\) and the signs of their literals.  A clause is \emph{ordinary}
if that subset is empty.  It is a \emph{test clause} if the subset
contains exactly one root and its literal is true at \(t\).  Let
\(O\) and \(Q\) be the ordinary and test indices among
\(\{1,\ldots,L\}\setminus\{i\}\).

Reveal the root type at every such index through \(L\).  At indices in
\(\{1,\ldots,L\}\setminus(\{i\}\cup O\cup Q)\), also reveal the
complete signed nonroot part.
Let \(\mathcal E\) be the sigma-field generated by this static
information.  Thus \(O,Q\) and the counts
\begin{equation}\label{eq:global-root-counts}
\begin{split}
 Z&=\#\{j\le L:j\ne i,\ C_j\text{ meets }R\},\\
 V&=\#\{j\le L:j\ne i,\ C_j\text{ meets at least two roots}\},
 \qquad d=\max\{1,Z\}
\end{split}
\end{equation}
are \(\mathcal E\)-measurable.  We call a clause counted by \(V\)
a \emph{collision}.  The ordinary clauses and the nonroot parts of
test clauses have not yet been revealed.

\begin{lemma}[Conditional product law]\label{lem:conditional-product}
Conditional on \(\mathcal E\), the clauses at indices in \(O\) are
independent uniform proper \(k\)-clauses on the \(u\) nonroots.
The nonroot parts at indices in \(Q\) are independent uniform proper
\((k-1)\)-clauses on the nonroots.  These two collections are mutually
independent.
\end{lemma}

\begin{proof}
For a prescribed root type with \(h\) roots and prescribed root signs,
there are exactly \(2^{k-h}\binom{u}{k-h}\) possible signed nonroot
parts.  Since complete proper \(k\)-clauses are equiprobable, these
parts are equiprobable conditional on the type.  Original clause
independence makes this conditional law a product over the indices.
Once the types are fixed, \(O\) and \(Q\) are fixed index sets.
Revealing the parts outside \(O\cup Q\) preserves the product law
on the other indices.  Their residual lengths are \(k\) on \(O\)
and \(k-1\) on \(Q\), as asserted.
\end{proof}

A test clause is already true under \(t\), so its nonroot part
imposes no condition on \(S_m\).
Consequently \(S_m\) is measurable with respect to
\[
 \mathcal G_m
 =\sigma\bigl(\mathcal E,\ C_j:j\in O,\ j\le m\bigr).
\]
We will bound the probability of \(A_m\) under this filtration, which
knows \(S_m\) but leaves the test parts hidden.  To bound the remaining
duration on \(A_m\), we also need to know whether \(A_m\) has occurred.
For that purpose, enlarge the filtration by revealing the past test parts:
\[
 \mathcal F_m^*
 =\sigma\bigl(\mathcal G_m,\ 
       \text{nonroot parts of }C_j:j\in Q,\ j\le m\bigr).
\]
The full baseline prefix \(B_m\), and thus \(A_m\), is
\(\mathcal F_m^*\)-measurable.  Under \(\mathcal G_m\), the past
test parts remain independent with their original conditional laws.
Under either filtration, the future ordinary clauses remain independent
uniform proper \(k\)-clauses.  Both assertions follow from
Lemma~\ref{lem:conditional-product}.  The static information includes
future nonordinary data, but it includes no future ordinary contents.

\subsection{Flipping a root gives an independent test}

Recall that \(q_d(S)\) is the probability that \(d\) independent
proper \((k-1)\)-clauses kill \(S\).  Here \(d=\max\{1,Z\}\) is already
known under \(\mathcal E\).  Put
\[
 \beta=\frac{k}{k-1},\qquad q_m=q_d(S_m).
\]
For \(j\in R\), let \(\mathcal I_{j,m}\) be the test indices through
\(m\) whose root is \(j\), and put \(s_{j,m}=|\mathcal I_{j,m}|\).
Then \(s_{j,m}\le Z\le d\).  Let \(J_m\subseteq R\) consist of roots
\(j\) for which no collision through \(m\) has exactly one root literal
true under \(t\), at root \(j\).  A collision can exclude at most one
root, so
\begin{equation}\label{eq:available-root-tests}
 |J_m|\ge k-V.
\end{equation}
The index sets and counts in this paragraph are
\(\mathcal E\)-measurable.

\begin{lemma}[Root-flip tests]\label{lem:root-flip-tests}
For \(i\le m<L\), on the event \(\{|J_m|\ge1\}\),
\begin{equation}\label{eq:root-flip-tests}
 \PP(A_m\mid\mathcal G_m)
 \le \mathbf1_{\{S_m\ne\varnothing\}}q_m^{|J_m|}.
\end{equation}
\end{lemma}

\begin{proof}
Suppose \(A_m\) occurs, and fix \(j\in J_m\).  The nonroot parts
at the indices in \(\mathcal I_{j,m}\) must kill \(S_m\).
If they did not, choose \(x\in S_m\) satisfying all of them.
Extend \(x\) by the root assignment \(t\), then flip root \(j\).
We check the clauses of \(B_m\).

A clause with no true root literal under \(t\) has a satisfied nonroot
part, because \((t,x)\) satisfies \(B_m\).  The flip does not change
that part.  A clause with a true literal on another root remains true.
Every remaining clause has \(j\) as its unique true root under \(t\).
It cannot be a collision, by \(j\in J_m\), so it is a test clause at
an index in \(\mathcal I_{j,m}\); its nonroot part is satisfied by
the choice of \(x\).  The flipped assignment therefore satisfies
\(B_m\) with root values different from \(t\), contradicting \(A_m\).

Conditional on \(\mathcal G_m\), the test parts belonging to distinct
roots are disjoint independent families of proper \((k-1)\)-clauses.
For nonempty \(S_m\), a family of \(s_{j,m}\le d\) tests kills it with
probability \(q_{s_{j,m}}(S_m)\le q_d(S_m)\).  This also covers
\(s_{j,m}=0\), whose killing probability is zero.  The necessary
tests for all roots in \(J_m\) thus have conditional probability at
most \(q_m^{|J_m|}\).  Finally \(A_m\) requires \(S_m\ne\varnothing\).
\end{proof}

The proof has kept the actual incidence count \(d\), rather than
discarding environments where it exceeds a chosen cutoff.  We next
bound duration with that same parameter; its moments will be averaged
only after all conditional estimates have been proved.

\subsection{From ordinary clauses to chronological duration}

For this comparison, keep the actual schedule and ordinary clauses
through \(L\).  At every index after \(L\), append an auxiliary
independent uniform proper \(k\)-clause on the \(u\) nonroots.
The continuation is independent of the entire original process.

For \(i\le m<L\) and \(S_m\ne\varnothing\), let \(K_m\) be the number
of ordinary clauses strictly after \(m\), including the first one
that completes the killing of the original set \(S_m\).  Use the
continuation if necessary.  Let \(U_m\) be the number of chronological
indices after \(m\) through that killing index.  Set \(K_m=U_m=0\)
if \(S_m=\varnothing\).  The conditional product law and the independent
continuation give
\begin{equation}\label{eq:ordinary-future-mean}
 \EE[K_m\mid\mathcal F_m^*]
 =\EE[K_m\mid\mathcal G_m]=\tau_k(S_m).
\end{equation}
There are at most \(Z\) nonordinary indices between \(m\) and the
kill: every one occurs before or at \(L\), and the deleted index
\(i\) is not after \(m\).  Thus, pathwise,
\begin{equation}\label{eq:chronological-domination}
 U_m\le K_m+Z,\qquad
 \EE[U_m\mid\mathcal F_m^*]\le\tau_k(S_m)+Z,\qquad
 \EE[U_m\mid\mathcal G_m]\le\tau_k(S_m)+Z.
\end{equation}
The continuation makes these statements valid even when no pre-cap
ordinary clause completes the kill.

\begin{lemma}[Weighted deletion delay]\label{lem:weighted-delay}
Let
\[
 W_m=\min\{L,\tau_k(S_m)+Z\}.
\]
This is \(\mathcal G_m\)-measurable.  For \(i\le m<L\), on \(A_m\),
\begin{equation}\label{eq:future-deletion-delay}
 \EE\left[\sum_{\ell=m+1}^{L-1}\mathbf1_{A_\ell}
       \,\middle|\,\mathcal F_m^*\right]\le W_m.
\end{equation}
Consequently
\begin{equation}\label{eq:conditional-weighted-delay}
 \EE[D_i^2\mid\mathcal E]
 \le3\sum_{m=i}^{L-1}
        \EE[\mathbf1_{A_m}W_m\mid\mathcal E].
\end{equation}
\end{lemma}

\begin{proof}
On \(A_m\), every later solution of \(B_\ell\) must have root values
\(t\) and a nonroot restriction in \(S_m\).  Once the future ordinary
clauses kill that original set, there can be no such solution, including
at the killing index.  If the kill is after \(L\), the remaining capped
interval ends first.  In either case the number of later deletion
events is at most \(U_m\), and is also at most \(L\).
Taking conditional expectations and using
\eqref{eq:chronological-domination} gives
\eqref{eq:future-deletion-delay}.

On \(A_m\), the set \(S_m\) is nonempty and therefore needs at least
one ordinary clause to be killed.  Hence \(\tau_k(S_m)\ge1\) and
\(W_m\ge1\).  Expand the square of \eqref{eq:delay-identity}.
For each cross term condition on the earlier \(\mathcal F_m^*\);
it contains \(A_m\).  The tower property gives
\begin{align*}
 \EE[D_i^2\mid\mathcal E]
 &=\sum_{m=i}^{L-1}\EE[\mathbf1_{A_m}\mid\mathcal E]\\
 &\quad+2\sum_{m=i}^{L-1}
  \EE\left[\mathbf1_{A_m}
   \EE\left[\sum_{\ell=m+1}^{L-1}\mathbf1_{A_\ell}
             \,\middle|\,\mathcal F_m^*\right]
             \,\middle|\,\mathcal E\right]\\
 &\le\sum_{m=i}^{L-1}
        \EE[\mathbf1_{A_m}(1+2W_m)\mid\mathcal E]\\
 &\le3\sum_{m=i}^{L-1}
        \EE[\mathbf1_{A_m}W_m\mid\mathcal E].\qedhere
\end{align*}
\end{proof}

Both the root-test probability and the remaining duration are now
controlled by the same set \(S_m\).  We still have to bound how long
that set can spend at each level of \(q_d(S_m)\), and then average the
incidence counts.
